% ===================================================================
% Paper M1 — Five axioms and their minimal geometry
% Base paper of the Master Equation Framework series.  v3 WORKING DRAFT, 10 Nov 2026
% (the fifth axiom added; Decision Log P19, P23, P25, P26).  Not canon.
% Compile: xelatex (three passes).  Certificates: M1_lean/*.lean, Lean 4.15.0 core.
% ===================================================================
\documentclass[11pt,a4paper]{article}
\usepackage{amsmath,amssymb,amsthm}
\usepackage[margin=25mm]{geometry}
\usepackage{enumitem}
\usepackage{booktabs}
\usepackage{fontspec}
\setmonofont{DejaVu Sans Mono}[Scale=0.78]
\usepackage{listings}
\lstset{basicstyle=\ttfamily\footnotesize,breaklines=true,columns=fullflexible,keepspaces=true,frame=single,framesep=4pt,xleftmargin=5pt,xrightmargin=5pt,extendedchars=true}
\usepackage[hidelinks]{hyperref}
%%=====================================================================
%% mef_macros.tex
%% Canonical macro set for the Master Equation Framework (MEF) series.
%% Version: 1.0 (April 2026)
%%
%% Single source of truth for symbol notation across MEF papers.
%% Derived from main.tex (thesis preamble) plus selected additions
%% from Paper XVI (number sets and rigour-classification marks).
%%
%% Usage: place %%=====================================================================
%% mef_macros.tex
%% Canonical macro set for the Master Equation Framework (MEF) series.
%% Version: 1.0 (April 2026)
%%
%% Single source of truth for symbol notation across MEF papers.
%% Derived from main.tex (thesis preamble) plus selected additions
%% from Paper XVI (number sets and rigour-classification marks).
%%
%% Usage: place %%=====================================================================
%% mef_macros.tex
%% Canonical macro set for the Master Equation Framework (MEF) series.
%% Version: 1.0 (April 2026)
%%
%% Single source of truth for symbol notation across MEF papers.
%% Derived from main.tex (thesis preamble) plus selected additions
%% from Paper XVI (number sets and rigour-classification marks).
%%
%% Usage: place \input{mef_macros.tex} in the preamble of every MEF
%% paper after the package loads but before \begin{document}.  Drift
%% between papers is prevented by editing only this file.
%%
%% Companion file: mef_field_content_preamble.tex (Section 2 reader-
%% orientation block).
%%=====================================================================

%% ---- Quantum Coherence Dilaton and its constituents ----
\newcommand{\QCD}{\Phi_{\mathrm{QC}}}           % Quantum Coherence Dilaton (canonical)
\newcommand{\consciousness}{\omega}             % consciousness coherence field
\newcommand{\zetaval}{12\pi}                    % quantum coherence coupling = chi(K_8) * pi

%% ---- Manifolds ----
\newcommand{\Mtwelve}{\mathcal{M}_{12}}         % 12D parent manifold
\newcommand{\Keight}{K_8}                       % 8D internal manifold
\newcommand{\CPtwo}{\mathbb{CP}^2}              % complex projective plane
\newcommand{\Stwo}{S^2}                         % 2-sphere
\newcommand{\Ttwo}{T^2}                         % 2-torus

%% ---- Discrete groups ----
\newcommand{\Ztwo}{\mathbb{Z}_2}                % pillowcase orbifold group
\newcommand{\Zfour}{\mathbb{Z}_4}

%% ---- Spin^c and bundles ----
\newcommand{\Spinc}{\text{Spin}^c}              % Spin^c structure
\newcommand{\Oline}[1]{\mathcal{O}(#1)}         % line bundle of degree #1

%% ---- Number fields (NB: \Q = mathbb Q is RATIONALS, not the dilaton; \QCD = \mathcal{Q} is the dilaton) ----
\newcommand{\C}{\mathbb{C}}
\newcommand{\R}{\mathbb{R}}
\newcommand{\Z}{\mathbb{Z}}
\newcommand{\Q}{\mathbb{Q}}

%% ---- Effective couplings, scales, moduli ----
\newcommand{\Veff}{V_{\text{eff}}}
\newcommand{\VCKM}{V_{\text{CKM}}}
\newcommand{\lambdaC}{\lambda_C}
\newcommand{\tautwo}{\tau_2}                    % stabilised aspect modulus
\newcommand{\thetaW}{\theta_W}                  % Weinberg angle
\newcommand{\Rw}{R_w}                           % S^2 radius (warped fibre)
\newcommand{\Rc}{R_c}                           % CP^2 radius (Calabi-like)
\newcommand{\Rfibre}{\mathcal{R}_{\text{fibre}}}
\newcommand{\Cfibre}{C_{\mathrm{fibre}}}
\newcommand{\Rsix}{\mathcal{R}_6}
\newcommand{\Fsq}{\mathcal{F}^2}
\newcommand{\centcorr}{\mathcal{C}_{32}}

%% ---- Spin^c holonomy phases ----
\newcommand{\Phiomega}{\phi_\consciousness}     % omega-axis holonomy
\newcommand{\Phipsi}{\phi_\psi}                 % psi-axis holonomy

%% ---- Arithmetic ----
\newcommand{\dodd}{d_{\mathrm{odd}}}            % odd divisor function

%% ---- Operators ----
\newcommand{\Tr}{\mathrm{Tr}}
\newcommand{\Vol}{\mathrm{Vol}}
\newcommand{\vol}{\mathrm{vol}}

%% ---- Rigour-classification marks ----
\newcommand{\rigmark}{\hfill\textbf{[Rigorous]}}
\newcommand{\dermark}{\hfill\textbf{[Derived]}}
\newcommand{\motmark}{\hfill\textbf{[Motivated]}}
\newcommand{\gapmark}{\hfill\textbf{[Gap]}}
\newcommand{\rigder}{\hfill\textbf{[Rigorous / Derived]}}
\newcommand{\dergap}{\hfill\textbf{[Derived / Gap]}}

%% ---- End of mef_macros.tex ----
 in the preamble of every MEF
%% paper after the package loads but before \begin{document}.  Drift
%% between papers is prevented by editing only this file.
%%
%% Companion file: mef_field_content_preamble.tex (Section 2 reader-
%% orientation block).
%%=====================================================================

%% ---- Quantum Coherence Dilaton and its constituents ----
\newcommand{\QCD}{\Phi_{\mathrm{QC}}}           % Quantum Coherence Dilaton (canonical)
\newcommand{\consciousness}{\omega}             % consciousness coherence field
\newcommand{\zetaval}{12\pi}                    % quantum coherence coupling = chi(K_8) * pi

%% ---- Manifolds ----
\newcommand{\Mtwelve}{\mathcal{M}_{12}}         % 12D parent manifold
\newcommand{\Keight}{K_8}                       % 8D internal manifold
\newcommand{\CPtwo}{\mathbb{CP}^2}              % complex projective plane
\newcommand{\Stwo}{S^2}                         % 2-sphere
\newcommand{\Ttwo}{T^2}                         % 2-torus

%% ---- Discrete groups ----
\newcommand{\Ztwo}{\mathbb{Z}_2}                % pillowcase orbifold group
\newcommand{\Zfour}{\mathbb{Z}_4}

%% ---- Spin^c and bundles ----
\newcommand{\Spinc}{\text{Spin}^c}              % Spin^c structure
\newcommand{\Oline}[1]{\mathcal{O}(#1)}         % line bundle of degree #1

%% ---- Number fields (NB: \Q = mathbb Q is RATIONALS, not the dilaton; \QCD = \mathcal{Q} is the dilaton) ----
\newcommand{\C}{\mathbb{C}}
\newcommand{\R}{\mathbb{R}}
\newcommand{\Z}{\mathbb{Z}}
\newcommand{\Q}{\mathbb{Q}}

%% ---- Effective couplings, scales, moduli ----
\newcommand{\Veff}{V_{\text{eff}}}
\newcommand{\VCKM}{V_{\text{CKM}}}
\newcommand{\lambdaC}{\lambda_C}
\newcommand{\tautwo}{\tau_2}                    % stabilised aspect modulus
\newcommand{\thetaW}{\theta_W}                  % Weinberg angle
\newcommand{\Rw}{R_w}                           % S^2 radius (warped fibre)
\newcommand{\Rc}{R_c}                           % CP^2 radius (Calabi-like)
\newcommand{\Rfibre}{\mathcal{R}_{\text{fibre}}}
\newcommand{\Cfibre}{C_{\mathrm{fibre}}}
\newcommand{\Rsix}{\mathcal{R}_6}
\newcommand{\Fsq}{\mathcal{F}^2}
\newcommand{\centcorr}{\mathcal{C}_{32}}

%% ---- Spin^c holonomy phases ----
\newcommand{\Phiomega}{\phi_\consciousness}     % omega-axis holonomy
\newcommand{\Phipsi}{\phi_\psi}                 % psi-axis holonomy

%% ---- Arithmetic ----
\newcommand{\dodd}{d_{\mathrm{odd}}}            % odd divisor function

%% ---- Operators ----
\newcommand{\Tr}{\mathrm{Tr}}
\newcommand{\Vol}{\mathrm{Vol}}
\newcommand{\vol}{\mathrm{vol}}

%% ---- Rigour-classification marks ----
\newcommand{\rigmark}{\hfill\textbf{[Rigorous]}}
\newcommand{\dermark}{\hfill\textbf{[Derived]}}
\newcommand{\motmark}{\hfill\textbf{[Motivated]}}
\newcommand{\gapmark}{\hfill\textbf{[Gap]}}
\newcommand{\rigder}{\hfill\textbf{[Rigorous / Derived]}}
\newcommand{\dergap}{\hfill\textbf{[Derived / Gap]}}

%% ---- End of mef_macros.tex ----
 in the preamble of every MEF
%% paper after the package loads but before \begin{document}.  Drift
%% between papers is prevented by editing only this file.
%%
%% Companion file: mef_field_content_preamble.tex (Section 2 reader-
%% orientation block).
%%=====================================================================

%% ---- Quantum Coherence Dilaton and its constituents ----
\newcommand{\QCD}{\Phi_{\mathrm{QC}}}           % Quantum Coherence Dilaton (canonical)
\newcommand{\consciousness}{\omega}             % consciousness coherence field
\newcommand{\zetaval}{12\pi}                    % quantum coherence coupling = chi(K_8) * pi

%% ---- Manifolds ----
\newcommand{\Mtwelve}{\mathcal{M}_{12}}         % 12D parent manifold
\newcommand{\Keight}{K_8}                       % 8D internal manifold
\newcommand{\CPtwo}{\mathbb{CP}^2}              % complex projective plane
\newcommand{\Stwo}{S^2}                         % 2-sphere
\newcommand{\Ttwo}{T^2}                         % 2-torus

%% ---- Discrete groups ----
\newcommand{\Ztwo}{\mathbb{Z}_2}                % pillowcase orbifold group
\newcommand{\Zfour}{\mathbb{Z}_4}

%% ---- Spin^c and bundles ----
\newcommand{\Spinc}{\text{Spin}^c}              % Spin^c structure
\newcommand{\Oline}[1]{\mathcal{O}(#1)}         % line bundle of degree #1

%% ---- Number fields (NB: \Q = mathbb Q is RATIONALS, not the dilaton; \QCD = \mathcal{Q} is the dilaton) ----
\newcommand{\C}{\mathbb{C}}
\newcommand{\R}{\mathbb{R}}
\newcommand{\Z}{\mathbb{Z}}
\newcommand{\Q}{\mathbb{Q}}

%% ---- Effective couplings, scales, moduli ----
\newcommand{\Veff}{V_{\text{eff}}}
\newcommand{\VCKM}{V_{\text{CKM}}}
\newcommand{\lambdaC}{\lambda_C}
\newcommand{\tautwo}{\tau_2}                    % stabilised aspect modulus
\newcommand{\thetaW}{\theta_W}                  % Weinberg angle
\newcommand{\Rw}{R_w}                           % S^2 radius (warped fibre)
\newcommand{\Rc}{R_c}                           % CP^2 radius (Calabi-like)
\newcommand{\Rfibre}{\mathcal{R}_{\text{fibre}}}
\newcommand{\Cfibre}{C_{\mathrm{fibre}}}
\newcommand{\Rsix}{\mathcal{R}_6}
\newcommand{\Fsq}{\mathcal{F}^2}
\newcommand{\centcorr}{\mathcal{C}_{32}}

%% ---- Spin^c holonomy phases ----
\newcommand{\Phiomega}{\phi_\consciousness}     % omega-axis holonomy
\newcommand{\Phipsi}{\phi_\psi}                 % psi-axis holonomy

%% ---- Arithmetic ----
\newcommand{\dodd}{d_{\mathrm{odd}}}            % odd divisor function

%% ---- Operators ----
\newcommand{\Tr}{\mathrm{Tr}}
\newcommand{\Vol}{\mathrm{Vol}}
\newcommand{\vol}{\mathrm{vol}}

%% ---- Rigour-classification marks ----
\newcommand{\rigmark}{\hfill\textbf{[Rigorous]}}
\newcommand{\dermark}{\hfill\textbf{[Derived]}}
\newcommand{\motmark}{\hfill\textbf{[Motivated]}}
\newcommand{\gapmark}{\hfill\textbf{[Gap]}}
\newcommand{\rigder}{\hfill\textbf{[Rigorous / Derived]}}
\newcommand{\dergap}{\hfill\textbf{[Derived / Gap]}}

%% ---- End of mef_macros.tex ----


\theoremstyle{plain}
\newtheorem{theorem}{Theorem}[section]
\newtheorem{proposition}[theorem]{Proposition}
\newtheorem{lemma}[theorem]{Lemma}
\newtheorem{corollary}[theorem]{Corollary}
\theoremstyle{definition}
\newtheorem{definition}[theorem]{Definition}
\newtheorem{axiom}{Axiom}
\theoremstyle{remark}
\newtheorem{remark}[theorem]{Remark}

\newcommand{\Hquat}{\mathbb{H}}
\newcommand{\ImH}{\operatorname{Im}\Hquat}
\newcommand{\Ad}{\operatorname{Ad}}
\newcommand{\Fix}{\operatorname{Fix}}
\newcommand{\Sol}{\mathrm{Sol}}
\newcommand{\rigR}{[\textup{R}]}
\newcommand{\rigD}{[\textup{D}]}
\newcommand{\rigM}{[\textup{M}]}
\newcommand{\rigG}{[\textup{G}]}

\title{Five axioms and their minimal geometry}
\author{Dhiren J. Master}
\date{10 November 2026 (v3, working draft)}

\begin{document}
\maketitle

\begin{abstract}
A unified theory must explain what is observed. We begin from five characteristics of the universe that need no instrument to recognise: it is aware of itself and closes on itself; it grows and decays; the order of its events matters; it never stands still; and what is seen of it is less than what is there. Taken as axioms and formalised as conditions on a compact Riemannian geometry $(T\times\Sigma)\times F$, with every point at which a behavioural word acquires a mathematical meaning isolated as a numbered definition, and with a lexicographic notion of minimality, they force a unique geometry. The time factor is the Sol torus bundle whose monodromy is the square of the Fibonacci matrix, so that the golden ratio is its growth rate; the internal fibre is the pillowcase orbifold with four cone points of angle $\pi$ carrying the subalgebras $\R,\C_i,\C_j,\C_k$ of the quaternions; the dimension is twelve; the base is $\CPtwo\times\Stwo$; and the geometry carries a $\Spinc$ structure with determinant line $\Oline{3}$. The constants $\pi$, $e$, $\phi$ and the primes two, three and five are outputs. As an application the minimal gauge manifold of Witten is recovered as a corollary. The principal consequence is that the twelve-dimensional manifold of the Master Equation programme is derived rather than postulated.
\end{abstract}

\section{What is observed, and the five axioms}\label{sec:intro}

There has been a long quest for a grand unified field theory. Whatever form such a theory takes, it must explain what we observe in the universe; that is the starting point, not a test applied at the end. We therefore begin with the observable characteristics of the universe in their simplest form, the ones that need no instrument and no training to recognise, and from them we state five axioms. From the five axioms, and a principle of choosing the simplest geometry that satisfies them, we derive in this paper a twelve-dimensional manifold; later papers of the series derive the field equation on it and its consequences. The purpose of this paper is the first step: to show that the geometry on which the programme rests is forced by what is observed, rather than postulated.

\subsection{What a child can observe}\label{sec:observe}

Five things are apparent to anyone.

I am aware of myself and of my surroundings, and I act accordingly. So does a tree, turning to the light; so does everything that lives. Whatever the universe is, it contains things that are aware of themselves, and so the universe as a whole is in some sense aware of itself: it closes on itself.

Everything changes, including the things we do not call alive: the weather, the rocks, the mountains, the continents, the earth itself. Nothing stays still.

There are cycles. Winter is followed by spring, summer, autumn and winter again; day follows night follows day. Yet no season is exactly like the last and no day repeats: the cycles return, but every cycle is different. Things grow and then fade, the seed to the tree to the fallen leaf, and the fading is as much a part of the cycle as the growing.

The order in which things are done matters, and a different order gives a different result. Anyone who has followed a recipe knows this: mix the ingredients and then cook them, and you have a cake; cook them and then mix, and you have something else. What has been done cannot be undone by doing it again in another order.

A thing and its shadow are not the same. A shadow on the ground tells a good deal about what casts it, but not everything: the shadow is flat and the thing is not, and the shadow does not say what the thing is made of. Two quite different things can cast the same shadow, and one thing casts very different shadows as the light moves round it. Whatever is seen is seen from somewhere, and it is less than what is there.

\subsection{The five axioms}\label{sec:fouraxioms}

We take these five observations as axioms, in the simplest words that carry them.

\begin{axiom}[closure]\label{ax:A1}
The universe closes on itself, so that it can be aware of itself.
\end{axiom}

\begin{axiom}[growth and decay]\label{ax:A2}
The universe grows and decays as a function of itself.
\end{axiom}

\begin{axiom}[order]\label{ax:A3}
The order of events matters, so that each instance is unique and cannot be foretold.
\end{axiom}

\begin{axiom}[drive]\label{ax:A4}
The universe is in constant change; it never stands still.
\end{axiom}

\begin{axiom}[shadow]\label{ax:A5}
We can see the shadow, but the shadow is missing information about what it is a shadow of.
\end{axiom}

Closure carries the cycles and the self-awareness; growth and decay carries the seed and the leaf; order carries the recipe; drive carries the weather and the mountains; shadow carries the fact that we see from where we stand, and see less than there is. That no cycle repeats exactly is not a further axiom but a consequence of the first and the fourth together, as Section~\ref{sec:drive} shows.

\subsection{Four numbers}\label{sec:numbers}

Four numbers will appear in the derivation, and it is worth saying at the outset where each comes from, because none of them is put in by hand.

The number $\pi$ is the ratio of a circle to its diameter. It enters through closure: a thing that closes on itself is a loop, and the simplest loop is a circle; and it enters a second time at the four special points of the internal geometry, which turn out to be cones whose angle is exactly half a full turn, $\pi$ (Section~\ref{sec:fibre}).

The number $e\approx2.718$ is the base of natural growth: a quantity that grows at a rate proportional to its own size grows as $e^{\kappa t}$. It enters through growth and decay, and nowhere else (Section~\ref{sec:A2}).

The golden ratio $\phi=(1+\sqrt5)/2\approx1.618$ is the number with $\phi^{2}=\phi+1$. It enters through the drive, as the rate at which the gentlest possible everlasting change stretches what it acts on (Section~\ref{sec:drive}).

The Fibonacci series $1,1,2,3,5,8,13,\dots$, in which each term is the sum of the two before it, grows at the rate $\phi$: the ratio of consecutive terms tends to $\phi$. It enters with the drive in a precise way. The drive, once the axioms have selected it, is the matrix $\bigl(\begin{smallmatrix}2&1\\1&1\end{smallmatrix}\bigr)$, and this matrix is the square of the Fibonacci matrix $\bigl(\begin{smallmatrix}1&1\\1&0\end{smallmatrix}\bigr)$, the matrix that sends a pair of consecutive Fibonacci numbers to the next pair. The least drive that never stands still is two steps of the Fibonacci recursion (Theorem~\ref{thm:fibonacci}).

\subsection{What this paper proves, and how it is organised}\label{sec:results}

The result is that the five axioms, together with a stated principle of minimality, force a unique geometry:
\begin{equation}\label{eq:M12}
\Mtwelve \;=\; (S^{1}_{\Sol}\times S^{3})\times\Keight, \qquad
\Keight \;=\; \bigl[(\CPtwo\times\Stwo)\times_{w}(\Ttwo/\Ztwo)\bigr]^{\Spinc}.
\end{equation}
In words: time is a circle along which a two-dimensional internal fibre is stretched and squeezed by the Fibonacci drive; the fibre, folded in half by the growth-and-decay exchange, is a pillowcase with four cone points carrying the four pieces of the quaternion algebra of the order axiom; the whole has twelve dimensions, every smaller number failing the order axiom; and the base over which the fibre sits is the sphere of the three imaginary quaternion units together with the space of their complex directions. Every discrete datum of~\eqref{eq:M12}, and the three primes two, three and five that appear in its arithmetic, are outputs.

Section~\ref{sec:axioms} turns the five sentences into mathematics, at a level intended to be followed by a first-year mathematics graduate, and isolates as numbered definitions the six points at which a behavioural word has to be given a mathematical meaning. Section~\ref{sec:minimal} defines minimality. Sections~\ref{sec:drive} to~\ref{sec:base} derive the time factor, the fibre, the dimension and the base; Section~\ref{sec:theorem} assembles the theorem; Section~\ref{sec:outlook} says what the geometry dictates for the physics. From Section~\ref{sec:drive} onward the register is that of a formal mathematics paper: every load-bearing step is proved or reduced to a cited classical theorem, the smallest case is worked by hand, and the finite and arithmetic parts are certified in Lean (Appendix~\ref{app:lean}). Rigour is tagged \rigR\ (proved, or a cited classical theorem), \rigD\ (derived under a stated reading) and \rigM\ (motivated).

\section{The five axioms as mathematics}\label{sec:axioms}

This section is written to be followed without specialist background. Each axiom is a sentence about behaviour; to use it we must say what the sentence means for a shape. We do this one axiom at a time, in words first and symbols second, and whenever a word has to be given a meaning that is a choice rather than a consequence, we say so and number the choice as a definition, so that a reader can see exactly where the axioms end and our decisions begin.

\subsection{The shapes we consider}

A \emph{geometry} here is a smooth shape, of some number of dimensions, on which lengths and angles can be measured, and which is allowed to have a few isolated sharp points (cone points, like the tip of a paper cone). We write it as a product of three pieces,
\begin{equation}\label{eq:family}
M \;=\; (T \times \Sigma) \times F ,
\end{equation}
where $T$ is a one-dimensional piece carrying time, $\Sigma$ is a piece carrying space, and $F$ is an internal piece that is not directly seen. That the internal piece is not directly seen is the fifth axiom at work before it has been stated: the shape is larger than what is seen of it, and Section~\ref{sec:A5} says what follows. A product of pieces means, as for a rectangle made of two intervals, that a point of $M$ is a point of each piece at once. Writing $M$ this way is only a way of naming the pieces: the way lengths are measured may tie the pieces together, and the second axiom will insist that it does.

\subsection{Closure}\label{sec:A1}

The sentence is: \emph{the universe closes on itself, so that it can be aware of itself}.

For a shape, ``closes on itself'' means two things. First, no piece goes on for ever and no piece has an edge: each is \emph{compact}, like a circle or a sphere rather than a line or a disc. Time, in particular, is a circle: there is no first moment and no last. Second, there is some loop in the shape that cannot be shrunk to a point, the way the loop round a doughnut's hole cannot be shrunk but a loop drawn on a ball can. A path that returns to its start and cannot be undone is the geometric form of ``aware of itself'': the shape refers back to itself in a way that is part of its structure, not an accident of how it is measured.

\begin{definition}[closure, formal]\label{def:A1}
A geometry satisfies A1 if $T$, $\Sigma$ and $F$ are compact, connected and without boundary, and the fundamental group of $M$ is non-trivial.
\end{definition}

The fundamental group is the collection of loops in $M$ up to continuous deformation; ``non-trivial'' means some loop does not deform to a point. We ask this of the whole shape, not of each piece: the time circle already supplies such a loop, and the space and internal pieces are free to have none of their own. On any compact shape there is always a shortest closed path \cite{LyusternikFet}; we ask for more, a loop that is topologically stuck, and that is the first of the choices this section records.

\subsection{Growth and decay}\label{sec:A2}

The sentence is: \emph{the universe grows and decays as a function of itself}.

``As a function of itself'' is the law of natural growth: a quantity $w$ whose rate of change is proportional to its own size, $\mathrm{d}w/\mathrm{d}s=\kappa w$, and so $w=w_{0}e^{\kappa s}$. This is the only place the number $e$ enters. ``Grows and decays'' says that both signs of $\kappa$ occur, and are two faces of one thing: there is a symmetry of the shape that swaps $\kappa$ for $-\kappa$. ``The universe'' grows: the growth is a property of the shape, not of something placed on it, which we take to mean that the growth law sits in the way lengths are measured, as a \emph{warp}: the size of one piece depends on where one is in another, the way the circles of latitude on a globe shrink towards the poles.

\begin{definition}[growth and decay, formal]\label{def:A2}
A geometry satisfies A2 if: (i) at least one product in~\eqref{eq:family} is a warped product, with a smooth function $w$ on one piece, $\nabla w\neq0$, entering the lengths measured on another; (ii) along the warping direction $\mathrm{d}w/\mathrm{d}s=\kappa\,w$ for a constant $\kappa\neq0$, so that $w=w_{0}e^{\kappa s}$; (iii) there is an involution $\iota$ of the warped piece with $\iota^{*}\kappa=-\kappa$.
\end{definition}

An involution is a symmetry that undoes itself when applied twice, like a reflection. Clause~(i) is the second of the choices: the growth is in the geometry rather than in a field placed on it. The involution $\iota$ is not identified here; Section~\ref{sec:fibre} will find that it is the fold that halves the internal piece and, on the algebra of the next axiom, the operation of conjugation.

\subsection{Order}\label{sec:A3}

The sentence is: \emph{the order of events matters, so that each instance is unique and cannot be foretold}.

Mathematics has an exact word for ``order matters'': \emph{non-commutative}. Two operations commute if doing them in either order gives the same result, as adding $2$ and adding $3$ do; they fail to commute if the order changes the outcome, as putting on socks and putting on shoes do. The axiom says that the operations the universe performs on its internal piece fail to commute. To make this precise we must say what those operations are.

\begin{definition}[the fibre algebra]\label{def:fibrealg}
The \emph{fibre algebra} $\mathcal{A}$ of a geometry is the collection of operations on the fibres of $F$, with the product of two operations being the result of doing one and then the other, in the order in which they occur in time.
\end{definition}

\begin{definition}[order, formal]\label{def:A3}
A geometry satisfies A3 if its fibre algebra is non-commutative: there are operations $a$ and $b$ with $ab\neq ba$.
\end{definition}

Two things follow at once and are not extra assumptions. Because $F$ sits over a compact piece, there are only finitely many independent operations: $\mathcal{A}$ is finite-dimensional. Because ``and then'' is how operations are combined, doing $a$ then $b$ then $c$ is the same whether one groups it as $(ab)c$ or $a(bc)$: $\mathcal{A}$ is associative. What the axiom does \emph{not} say is that every operation can be undone; that it can will be a consequence of closure (Section~\ref{sec:fibre}), and it is that consequence, together with a classical theorem of Frobenius \cite{Frobenius1878}, that will force $\mathcal{A}$ to be the quaternions, the four-dimensional number system with units $1,i,j,k$ and $i^{2}=j^{2}=k^{2}=ijk=-1$, in which $ij=k$ but $ji=-k$. The second half of the sentence is a consequence of the first: for two observables that fail to commute, the uncertainties in measuring them together are bounded below by half the size of their commutator \cite{Robertson1929}, so ``cannot be foretold'', in the quantum sense, is a theorem of ``order matters''. Unpredictability in the other sense, of a system whose future is sensitive to its present, belongs to the fourth axiom.

Two further definitions say \emph{where} the fibre algebra acts, and these are the third and fourth choices. A shape has at every point its tangent directions, and, if it is the right kind of shape, also its spinors, the objects that carry the half-integer spin of a fermion. The order axiom will be read as saying that the non-commutative algebra is the algebra of the shape's own tangent directions and spinors, not an abstract bookkeeping device.

\begin{definition}[the algebra acts on spinors and on tangent vectors]\label{def:levels}
The fibre algebra acts on the geometry at two levels: at the spinor level, as the algebra generated by the Clifford action on a spinor bundle over $M$; and at the tangent level, as an almost-quaternionic structure on $M$, a rank-three subbundle of $\operatorname{End}(TM)$ locally spanned by anticommuting almost-complex structures $I,J,K$ with $IJ=K$ \cite{Salamon1982}, whose restriction to $F$ carries the full algebra.
\end{definition}

An almost-complex structure $I$ is a rule that rotates every tangent direction by a quarter-turn, $I^{2}=-1$, the way multiplication by $i$ rotates the complex plane; a quaternionic structure is three such rules $I,J,K$ obeying the quaternion relations. The reader need only hold on to the picture: the shape carries, at every point, a copy of the quaternion multiplication table acting on its directions.

\begin{definition}[orbits are realised]\label{def:orbits}
A homogeneous space $G/H$ formed from the symmetry group $G$ of the fibre algebra and the stabiliser $H$ of a structure selected on it is \emph{realised} by the geometry: it appears as a piece of the base over which $F$ is warped.
\end{definition}

The picture here is simpler than the words. The quaternions have a symmetry group, the rotations of the three imaginary units $i,j,k$ among themselves. If something in the geometry picks out one unit, say $i$, then the set of all units it \emph{could} have picked is a sphere, and the definition says that this sphere is part of the shape, not merely a list of alternatives. That is the fourth choice, and it is what produces the two-sphere and the complex projective plane of Section~\ref{sec:base}, where the other possible reading is set out.

\subsection{Drive}\label{sec:A4}

The sentence is: \emph{the universe is in constant change; it never stands still}.

For a shape with a time circle this means a flow: a rule moving every point along as time advances. ``Never stands still'' means the flow has no fixed points. ``Constant change'' means more than motion: a planet in a perfectly regular orbit moves but changes nothing, and we take the axiom to exclude that. The mathematical measure of genuine change is \emph{entropy} of the flow: it is positive when two points that start close are driven apart, so that the far future cannot be read off from the present. And ``never comes to the same point again'' means that although time itself is a circle, a typical state of the universe does not return: the points that do return are exceptional, of measure zero.

\begin{definition}[drive, formal]\label{def:A4}
A geometry satisfies A4 if $T\times F$ carries a smooth flow $\varphi_{t}$ with no fixed points, positive topological entropy, and a set of periodic points of measure zero.
\end{definition}

The flow is placed on time together with the internal piece, not on time alone, because a one-dimensional circle cannot carry a flow that drives points apart; the drive is what time does to the internal piece as it goes round. This is the first point at which the order axiom and the drive act on the same object. The strong form of the drive that the construction uses, a flow that stretches in one direction and squeezes in another, is not assumed; Section~\ref{sec:drive} derives it, and shows that the gentlest such flow is the Fibonacci drive of Section~\ref{sec:numbers}.

\subsection{Shadow}\label{sec:A5}

The sentence is: \emph{we can see the shadow, but the shadow is missing information about what it is a shadow of}.

For a shape, the sentence says three things. First, there is a place we see from. We are part of the shape; we see along the pieces that carry time and space, and the internal piece is not seen directly. We also persist: we are here after each circuit of time, so the place we see from is a point of the internal piece that the drive returns to itself. Second, what reaches us of the internal piece is its shadow: what is left of a thing in the whole shape when its position in the internal piece is forgotten. Forgetting loses information, so two different things in the whole shape can leave the same shadow, and one thing can leave different shadows. Third, the other way about: whatever we see is the shadow of something in the whole shape, and the whole shape is regular even where the shadow looks broken.

\begin{definition}[shadow, formal]\label{def:A5}
A geometry satisfies A5 if: (i) there is a point $o$ of the fibre of $F$, returned to itself by the drive of A4 after each circuit of $T$, from which observation is made, and what is observed is the image of what is on $M$ under the projection $\pi\colon M\to T\times\Sigma$ that forgets the internal piece; (ii) the growth law of A2 lies on the piece of $F$ that carries the involution of Definition~\ref{def:A2}(iii), and on no other: no length measured on $T\times\Sigma$, or on the base over which that piece sits, depends on position in it; (iii) every thing observed, including what appears singular in $T\times\Sigma$, is the image under $\pi$ of a thing on $M$ that is regular.
\end{definition}

The projection $\pi$ is not one-to-one, since the internal piece is not a point: that is the loss of information the sentence speaks of, and it is a consequence of clause~(i), not a further condition. Clause~(iii) says that something regular lies above every shadow; it does not say that it is the only thing that could. In this paper the shape is fixed by the five axioms and the principle of minimality, and not by its shadow.

Clauses~(i) and~(ii) are the fifth and sixth of the choices this section records. Clause~(i) says where we stand, and the reason for it is in the sentence itself: a shadow is seen from somewhere, and an observer who is still there after each circuit must sit where the drive brings it back. Section~\ref{sec:drive} finds that the drive returns exactly one point of the fibre to itself, so the place is not a choice among several. Clause~(ii) says where the growth law lies. Definition~\ref{def:A2} asks for a warp on one piece entering the lengths of another and does not say which; it already places the involution on ``the warped piece'', and Section~\ref{sec:fibre} finds that the piece carrying the involution is the torus that folds to the pillowcase. The fifth axiom adds that the warp is there and nowhere else. The reason is physical, and we give it in words since this paper has no forces in it: gravity is the curvature of the whole twelve-dimensional shape and reaches everywhere in it, which is why it is weak; the other forces are carried by the base and are local to it, and the growth law does not touch them. \rigD

\subsection{What has been chosen}\label{sec:scope}

The five sentences have been turned into five conditions, and in doing so six choices were made and numbered: that self-reference means a loop that cannot be shrunk (Definition~\ref{def:A1}); that growth is in the geometry (Definition~\ref{def:A2}(i)); that the non-commutative algebra is the algebra of the shape's own spinors and tangent directions (Definition~\ref{def:levels}); that the sphere of alternatives is part of the shape (Definition~\ref{def:orbits}); that we see from a point the drive returns to itself (Definition~\ref{def:A5}(i)); and that the growth law lies on the folded piece and nowhere else (Definition~\ref{def:A5}(ii)). Every result below that depends on one of these says so. Nothing in the five conditions mentions a gauge group, a particle, a generation or a coupling constant, and nothing names a factor of the manifold that will be derived. One quantity is not constrained by any axiom: the overall size of the shape. The definition of minimality in the next section is made independent of size, and the final theorem states the size as the one input the axioms do not supply.

\section{Minimality}\label{sec:minimal}

\begin{definition}[modulus]\label{def:modulus}
A \emph{modulus} of a geometry satisfying A1--A4 is a continuous parameter of its structure data (the warping function up to~\eqref{def:A2}(ii), the complex and $\Spinc$ structures, the orbifold data, the flux integers) that survives the stationarity conditions of the geometry's energy functional, taken modulo diffeomorphism and overall scale.
\end{definition}

\begin{definition}[minimal]\label{def:minimal}
Among geometries satisfying A1--A4, $M$ is \emph{minimal} if it is selected by the following criteria applied in order, each used only to break ties left by the previous ones: (1) parametric: $M$ has no modulus; (2) dynamical: the drive of $M$ has the least positive topological entropy; (3) arithmetic: the topological integers of $F$ (its Euler characteristic, its flux quanta, the coefficients of the polynomial fixing its aspect ratio) lie in $\{2^{a}3^{b}\}$, the prime five entering only through the drive; (4) dimensional: $M$ has least total dimension, and then least $\lvert\pi_{1}(\Sigma)\rvert$.
\end{definition}

Criterion~(1) is the statement that a geometry with a dial would admit two universes differing only by that dial. Criterion~(2) selects the gentlest drive that still drives; it is of the same character as the last clause of~(4), the smallest fundamental group, and both express the instinct that the geometry carries no unforced structure. Criterion~(3) is a filter whose content is shown in Section~\ref{sec:theorem} to be an output of the selection rather than an input. Overall scale is exempt by construction.

\section{The drive}\label{sec:drive}

\subsection{Hyperbolicity}

Let the fibre of $F$ carrying the drive be a torus $\Ttwo=\R^{2}/\Z^{2}$; that it is one is proved in Section~\ref{sec:fibre} (Proposition~\ref{prop:surface}), and nothing in the present section is used there. The drive of A4 acts on $\Ttwo$ by its return map, an element $M\in\mathrm{SL}_{2}(\Z)$ with trace $t$.

\begin{proposition}[A4 is hyperbolicity]\label{prop:hyp}
A4 holds for the return map $M$ if and only if $\lvert t\rvert>2$. \rigR
\end{proposition}

\begin{proof}
The characteristic polynomial is $\lambda^{2}-t\lambda+1$. If $\lvert t\rvert<2$ the eigenvalues are complex of modulus one and $M$ has finite order: every point is periodic and the entropy vanishes, contradicting two clauses of Definition~\ref{def:A4}. If $\lvert t\rvert=2$ the eigenvalues are $\pm1$ and $M$ is a shear: orbits separate linearly, the entropy vanishes. If $\lvert t\rvert>2$ the eigenvalues are real, $\lambda$ and $\lambda^{-1}$ with $\lvert\lambda\rvert>1$; $M$ is Anosov, its entropy is $\log\lvert\lambda\rvert>0$, its periodic points are exactly the rational points, which have measure zero, and the induced flow on the mapping torus has no fixed points \cite{KatokHasselblatt}.
\end{proof}

\subsection{The least drive and its uniqueness}

Criterion~(2) of Definition~\ref{def:minimal} selects the least entropy, hence the least $\lvert t\rvert>2$, which is $t=3$ (certificate \texttt{L1\_trace3\_min}). The content of the section is that this selection is unique.

The selection is not that trace three is the only hyperbolic trace, which is false, but that it is the least, and that at the least trace there is nothing to choose: the theorem says that the gentlest drive which still drives is a single object up to relabelling the torus. The two proofs are independent in kind, one through ideal classes and one through reduced forms, and they are followed by the explicit relabellings that show the uniqueness at work.

\begin{theorem}[the drive]\label{thm:drive}
Every $M\in\mathrm{SL}_{2}(\Z)$ of trace three is conjugate in $\mathrm{SL}_{2}(\Z)$ to $M_{0}=\bigl(\begin{smallmatrix}2&1\\1&1\end{smallmatrix}\bigr)$. Consequently the minimal drive is unique; its eigenvalues are $\phi^{\pm2}$ with $\phi=(1+\sqrt5)/2$, its entropy is $2\log\phi$, its splitting field is $\Q(\sqrt5)$, and the time factor of a minimal geometry is the Sol torus bundle $S^{1}_{\Sol}$ with this monodromy. \rigR
\end{theorem}

\begin{proof}
Two routes. First, by the theorem of Latimer and MacDuffee \cite{LatimerMacDuffee1933} as sharpened by Taussky \cite{Taussky1949}, the $\mathrm{SL}_{2}(\Z)$-conjugacy classes of integer matrices with characteristic polynomial $p$ are in bijection with the narrow ideal classes of the order $\Z[\lambda]$, $\lambda$ a root of $p$. Here $p=\lambda^{2}-3\lambda+1$ with root $\phi^{2}=\phi+1$, so $\Z[\phi^{2}]=\Z[\phi]$, the ring of integers of $\Q(\sqrt5)$, which has class number one; its fundamental unit $\phi$ has norm $-1$, so the narrow class number is also one. There is a single class.

Second, by Gauss \cite{Gauss1801} (modern statement in \cite{Sarnak1982}), $M=\bigl(\begin{smallmatrix}a&b\\c&d\end{smallmatrix}\bigr)\mapsto cx^{2}+(d-a)xy-by^{2}$ carries conjugacy classes of trace $t$ bijectively to equivalence classes of binary quadratic forms of discriminant $t^{2}-4$. At $t=3$ the discriminant is five, squarefree, so all forms are primitive, and a form is reduced when $0<b<\sqrt5$ and $\sqrt5-b<2\lvert a\rvert<\sqrt5+b$. Over the integers these inequalities force $(a,b,c)\in\{(1,1,-1),(-1,1,1)\}$ (certificate \texttt{L1\_reduced\_disc5}), the two forms lie on one Gauss cycle (\texttt{L1\_cycle\_disc5}), and classes correspond to cycles. Again a single class.

Smallest case by hand. The candidates $M_{0}$ and $M_{0}^{-1}=\bigl(\begin{smallmatrix}1&-1\\-1&2\end{smallmatrix}\bigr)$ are conjugate by the quarter-turn $S=\bigl(\begin{smallmatrix}0&1\\-1&0\end{smallmatrix}\bigr)$: $SM_{0}S^{-1}=M_{0}^{-1}$ (\texttt{L1\_conj\_M0\_inv}); $M_{0}$ is symmetric, so its transpose adds nothing. The eigenvalue and field statements are the characteristic equation; the M\"obius action $\tau\mapsto(2\tau+1)/(\tau+1)$ has fixed points solving $\tau^{2}-\tau-1=0$, namely $\phi$ and $-1/\phi$.
\end{proof}

\subsection{Worked cases}\label{sec:drive-worked}

\emph{Three trace-three matrices and their conjugators.} The theorem asserts that any $A\in\mathrm{SL}_{2}(\Z)$ of trace three is $SM_{0}S^{-1}$ for some $S$. Three instances, each checked by multiplying out:
\begin{align*}
A&=\begin{pmatrix}3&-1\\1&0\end{pmatrix},& S&=\begin{pmatrix}-2&1\\-1&0\end{pmatrix};\\
A&=\begin{pmatrix}1&-1\\-1&2\end{pmatrix},& S&=\begin{pmatrix}-1&-1\\2&1\end{pmatrix};\\
A&=\begin{pmatrix}0&-1\\1&3\end{pmatrix},& S&=\begin{pmatrix}-2&-1\\3&1\end{pmatrix}.
\end{align*}
In each case $\det S=1$ and $SAS^{-1}=M_{0}$; the second instance is $M_{0}^{-1}$ with a conjugator different from the quarter-turn of the proof, showing that the conjugator is not unique while the class is.

\emph{The reduction of a form.} The form of $A=\bigl(\begin{smallmatrix}3&-1\\1&0\end{smallmatrix}\bigr)$ is $cx^{2}+(d-a)xy-by^{2}=x^{2}-3xy+y^{2}$, discriminant $9-4=5$. It is not reduced ($b=-3<0$). The Gauss step $\rho(a,b,c)=(c,b',c')$ with $b'\equiv-b\pmod{2c}$ and $b'$ in the reduced window gives $(1,-3,1)\mapsto(1,1,-1)$, one of the two reduced forms; the next step gives $(-1,1,1)$ and the next returns to $(1,1,-1)$. The cycle has length two and contains both reduced forms, so there is one class.

\emph{The eigenvalues and the entropy.} $\lambda^{2}-3\lambda+1=0$ gives $\lambda=(3\pm\sqrt5)/2$; since $\phi^{2}=\phi+1=(3+\sqrt5)/2$, the eigenvalues are $\phi^{2}\approx2.618$ and $\phi^{-2}\approx0.382$, and the entropy is $\log\phi^{2}=2\log\phi\approx0.962$. For comparison the next admissible trace, four, has eigenvalue $2+\sqrt3\approx3.732$ and entropy $\approx1.317$: larger, as criterion~(2) requires it to be, and with splitting field $\Q(\sqrt3)$ rather than $\Q(\sqrt5)$, so that the prime five is tied to the minimal choice.

\emph{The half-lattice evaluations.} With $M_{0}=\bigl(\begin{smallmatrix}2&1\\1&1\end{smallmatrix}\bigr)$: $(0,0)\mapsto(0,0)$; $(\tfrac12,0)\mapsto(1,\tfrac12)\equiv(0,\tfrac12)$; $(0,\tfrac12)\mapsto(\tfrac12,\tfrac12)$; $(\tfrac12,\tfrac12)\mapsto(\tfrac32,1)\equiv(\tfrac12,0)$. Reduced mod two, $\bar M_{0}=\bigl(\begin{smallmatrix}0&1\\1&1\end{smallmatrix}\bigr)$, $\bar M_{0}^{2}=\bigl(\begin{smallmatrix}1&1\\1&0\end{smallmatrix}\bigr)$, $\bar M_{0}^{3}=I$. The even-trace matrix $\bigl(\begin{smallmatrix}2&1\\3&2\end{smallmatrix}\bigr)$ reduces to $\bigl(\begin{smallmatrix}0&1\\1&0\end{smallmatrix}\bigr)$, a transposition of order two: it swaps two corners and fixes two, and cannot cycle three.


\begin{theorem}[the Fibonacci drive]\label{thm:fibonacci}
Let $F=\bigl(\begin{smallmatrix}1&1\\1&0\end{smallmatrix}\bigr)$ be the Fibonacci matrix, so that $F\binom{F_{n}}{F_{n-1}}=\binom{F_{n+1}}{F_{n}}$ and $F^{n}=\bigl(\begin{smallmatrix}F_{n+1}&F_{n}\\F_{n}&F_{n-1}\end{smallmatrix}\bigr)$. Then $M_{0}=F^{2}$. The minimal drive is two steps of the Fibonacci recursion; its eigenvalue $\phi^{2}$ is the square of the Fibonacci growth rate $\phi$, and $F$ itself, having determinant $-1$ and trace $1$, is not an element of $\mathrm{SL}_{2}(\Z)$: $M_{0}$ is the least power of $F$ that is. \rigR
\end{theorem}

\begin{proof}
$F^{2}=\bigl(\begin{smallmatrix}1\cdot1+1\cdot1&1\cdot1+1\cdot0\\1\cdot1+0\cdot1&1\cdot1+0\cdot0\end{smallmatrix}\bigr)=\bigl(\begin{smallmatrix}2&1\\1&1\end{smallmatrix}\bigr)$ (certificate \texttt{L1\_M0\_is\_Fib\_squared}). The entry recursion $(XF)_{11}=X_{11}+X_{12}$ is the Fibonacci law (\texttt{L1\_Fib\_recursion}), and $F^{2},\dots,F^{6}$ carry $(2,1,1,1)$, $(3,2,2,1)$, $(5,3,3,2)$, $(8,5,5,3)$, $(13,8,8,5)$ (\texttt{L1\_Fib\_powers}). The characteristic polynomial of $F$ is $\lambda^{2}-\lambda-1$ with roots $\phi,-1/\phi$; squaring gives $\phi^{2},\phi^{-2}$, the eigenvalues of Theorem~\ref{thm:drive}. $\det F=-1$, $\operatorname{tr}F=1$, $\det F^{2}=1$, $\operatorname{tr}F^{2}=3$ (\texttt{L1\_Fib\_not\_SL2}).
\end{proof}

The theorem says in what sense the golden ratio and the Fibonacci series are forced rather than found: the least everlasting change that the axioms allow is the Fibonacci recursion applied twice, and $\phi$ is its growth rate. That the cat map is the square of the Fibonacci matrix is classical; what is new is that the axioms select it.

\begin{proposition}[the corner three-cycle]\label{prop:threecycle}
For $M\in\mathrm{SL}_{2}(\Z)$, the reduction $\bar M\in\mathrm{GL}_{2}(\mathbb{F}_{2})\cong S_{3}$ has order three if and only if $t$ is odd. For $M_{0}$, $\bar M_{0}=\bigl(\begin{smallmatrix}0&1\\1&1\end{smallmatrix}\bigr)$ has order three; on the half-lattice points $P_{1}=(0,0)$, $P_{2}=(\tfrac12,0)$, $P_{3}=(0,\tfrac12)$, $P_{4}=(\tfrac12,\tfrac12)$ the map $M_{0}$ fixes $P_{1}$ and permutes the others as the cycle $(P_{2}\,P_{3}\,P_{4})$. \rigR
\end{proposition}

\begin{proof}
Over $\mathbb{F}_{2}$ the characteristic polynomial is $\lambda^{2}+\bar t\lambda+1$: for even $t$ it is $(\lambda+1)^{2}$ and $\bar M$ has order one or two; for odd $t$ it is $\lambda^{2}+\lambda+1$, irreducible, with roots of order three in $\mathbb{F}_{4}^{\times}$. The finite check over all of $\mathrm{SL}_{2}(\mathbb{F}_{2})$ is \texttt{L1\_odd\_trace\_iff\_order\_three}. The action on the half-lattice: $M_{0}P_{1}=P_{1}$; $M_{0}(\tfrac12,0)=(1,\tfrac12)\equiv P_{3}$; $M_{0}(0,\tfrac12)=(\tfrac12,\tfrac12)=P_{4}$; $M_{0}(\tfrac12,\tfrac12)=(\tfrac32,1)\equiv P_{2}$.
\end{proof}

The drive returns no other point of the torus to itself. A point $x$ with $M_{0}x\equiv x$ satisfies $(M_{0}-I)x\in\Z^{2}$, and $\det(M_{0}-I)=(2-1)(1-1)-1\cdot1=-1$, so $M_{0}-I$ is invertible over the integers and $x\in\Z^{2}$: the only such point is $P_{1}$ (\texttt{L1\_one\_fixed\_point}). This is the place we see from (Definition~\ref{def:A5}(i)). \rigR

The proposition is used in Section~\ref{sec:fibre} to assign the corner algebras. We note that the orientation of the three-cycle, $P_{2}\to P_{3}\to P_{4}$ against $P_{2}\to P_{4}\to P_{3}$, is the difference between $M_{0}$ and $M_{0}^{-1}$, which Theorem~\ref{thm:drive} shows to be a single conjugacy class: the orientation is a convention with no invariant content (Appendix~\ref{app:labels}).

\section{The fibre}\label{sec:fibre}

\subsection{The fibre algebra is the quaternions}

\begin{lemma}[the norm form]\label{lem:norm}
Let $\iota$ act on the fibre algebra $\mathcal{A}$ as an anti-involution $a\mapsto a^{*}$ (the algebraic image of exchanging growth and decay: composition run backwards). Then $N(a)=aa^{*}$ takes values in the fixed set of $*$, and is positive-definite if and only if the unit group of every corner subalgebra of the form $\R\oplus\R u$, $u^{*}=-u$, is compact. \rigR
\end{lemma}

\begin{proof}
$(aa^{*})^{*}=aa^{*}$, so $N$ is $*$-fixed. On $\R\oplus\R u$ the unit group is $\{x+yu:\,x^{2}-u^{2}y^{2}=1\}$, a circle if $u^{2}=-1$ and a hyperbola if $u^{2}=+1$; compactness is $u^{2}=-1$, i.e.\ $N(u)=1>0$ on every such line, which is positive-definiteness of $N$ on the span.
\end{proof}

The proposition that follows is the place where the first axiom does algebraic work. Non-commutativity alone admits the split quaternions, in which two of the three imaginary units square to $+1$ and whose unit groups are hyperbolas; what rules them out is that a corner whose unit group is a hyperbola is not compact, and A1 requires every factor, the corner fibres included, to close.

\begin{proposition}[$\mathcal{A}=\Hquat$]\label{prop:H}
Under A1, A2 and A3 the fibre algebra is the quaternions. \rigD
\end{proposition}

\begin{proof}
By Definition~\ref{def:fibrealg} $\mathcal{A}$ is finite-dimensional and associative. A1 requires every factor, the corner fibres included, to be compact, so by Lemma~\ref{lem:norm} the norm form of A2's anti-involution is positive-definite and $a^{-1}=a^{*}/N(a)$ for $a\neq0$: $\mathcal{A}$ is a division algebra. By Frobenius \cite{Frobenius1878} a finite-dimensional associative real division algebra is $\R$, $\C$ or $\Hquat$, and A3 excludes the two commutative cases. The alternative that compactness excludes is the split algebra $M_{2}(\R)$, which carries the same $\pm1,\pm i,\pm j,\pm k$ pattern with non-compact unit groups on two of its three imaginary lines. The relations $i^{2}=j^{2}=k^{2}=ijk=-1$ and the non-commutativity $ij\neq ji$ are certified in an integer model (\texttt{L2\_hamilton}, \texttt{L2\_noncomm}), as is positive-definiteness of $N(q)=q\bar q$ (\texttt{L2\_norm\_posdef}).
\end{proof}

\subsection{The four corners}

\begin{proposition}[corner count]\label{prop:corners}
The subalgebras of $\Hquat$ generated by a single unit of $Q_{8}=\{\pm1,\pm i,\pm j,\pm k\}$ are $\R$, $\C_{i}$, $\C_{j}$, $\C_{k}$, in bijection with $Q_{8}/\{\pm1\}\cong\Z_{2}\times\Z_{2}$: there are exactly four. The group of inner automorphisms by units, $V=\{\mathrm{id},\Ad_{i},\Ad_{j},\Ad_{k}\}$ with $\Ad_{u}(q)=uqu^{-1}$, is the same Klein four-group. \rigR
\end{proposition}

\begin{proof}
$u$ and $-u$ generate the same line; four lines, four classes (\texttt{L2\_corner\_count}). $\Ad_{u}=\Ad_{-u}$ and $\Ad_{i}\Ad_{j}=\Ad_{k}$ by Hamilton's relations.
\end{proof}

\begin{proposition}[fixed subalgebras]\label{prop:fixed}
Let $\iota$ act on $\Hquat$ as conjugation $q\mapsto\bar q$. Among the eight elements of $V\times\langle\iota\rangle$, those whose fixed set in $\Hquat$ is a commutative subalgebra are exactly $\iota$, $\Ad_{i}$, $\Ad_{j}$, $\Ad_{k}$, with fixed sets $\R$, $\C_{i}$, $\C_{j}$, $\C_{k}$ respectively; $\mathrm{id}$ fixes all of $\Hquat$, which is not commutative, and $\iota\circ\Ad_{u}$ fixes the three-dimensional $u^{\perp}$, which is not closed under multiplication. \rigR
\end{proposition}

\begin{proof}
Direct computation on $q=a+bi+cj+dk$: $\bar q=q$ iff $b=c=d=0$; $\Ad_{i}q=a+bi-cj-dk=q$ iff $c=d=0$, and cyclically; $\overline{\Ad_{i}q}=q$ iff $b=0$, and $jk=i\notin i^{\perp}$. Certified in \texttt{L2\_fix\_conj}, \texttt{L2\_fix\_Ad\_i}, \texttt{L2\_fix\_Ad\_j}, \texttt{L2\_fix\_Ad\_k}, \texttt{L2\_fix\_conj\_Ad\_i}, \texttt{L2\_iperp\_not\_closed}.
\end{proof}

\subsection{The fibre is a surface, and a torus}

\begin{proposition}[the fibre]\label{prop:surface}
Let the warped fibre $X$ of A2 be compact and connected, with the involution $\iota$ having exactly four isolated fixed points and the drive of A4 acting as an Anosov diffeomorphism commuting with $\iota$. Then the least-dimensional such $X$ is the torus $\Ttwo$, on which $\iota$ is inversion $z\mapsto-z$ and the drive is a hyperbolic element of $\mathrm{SL}_{2}(\Z)$. \rigD
\end{proposition}

\begin{proof}
If $\dim X=2$, Franks' theorem \cite{Franks1970} states that every Anosov diffeomorphism of a compact surface lives on $\Ttwo$ and is topologically conjugate to a hyperbolic linear automorphism. An involution of $T^{n}$ with only isolated fixed points is conjugate to an affine map $x\mapsto Ax+c$ with $A^{2}=I$ and no eigenvalue one, i.e.\ $A=-I$, with $\lvert\det(I-A)\rvert=2^{n}$ fixed points by the Lefschetz formula; four fixed points force $n=2$ (\texttt{L2\_two\_torsion}) and $\iota$ is inversion. On a three-torus the count is eight, excluding $\dim X=3$ \cite{FranksNewhouse}. Criterion~(4) of Definition~\ref{def:minimal} selects the least dimension; the selection is unconditional, resting on Franks' theorem at the minimal case, while the non-existence of exotic Anosov manifolds in dimension four and above, an open conjecture, is not needed.
\end{proof}

That $\iota$ has exactly four isolated fixed points is Proposition~\ref{prop:corners} read geometrically: the fixed points of $\iota$ are the points at which the algebra reduces, and the next proposition makes the reduction precise.

\subsection{The collapse at the corners}

\begin{proposition}[the corner algebras]\label{prop:collapse}
Let $\langle\iota,M_{0}\rangle$ act on $\Ttwo$ with $\iota$ inversion and $M_{0}$ the drive of Theorem~\ref{thm:drive}, and let the algebra at a point be the subalgebra of $\Hquat$ fixed by the stabiliser of that point, $\iota$ acting as conjugation and the drive as the order-three inner automorphism $\Ad_{w}$, $w=\tfrac12(1+i+j+k)$. Then the algebra is $\R$ at $P_{1}$, $\C_{i}$, $\C_{j}$, $\C_{k}$ at the three other fixed points of $\iota$, permuted cyclically by the drive, and $\Hquat$ at every point with trivial stabiliser. \rigD
\end{proposition}

\begin{proof}
The principle is A4 read on A3: order has nothing to act on where nothing moves, and the surviving algebra at a point is what its stabiliser fixes. $P_{1}$ is fixed by $\iota$ and by $M_{0}$ (Proposition~\ref{prop:threecycle}); the subalgebra fixed by conjugation is $\R$ (Proposition~\ref{prop:fixed}), and $\R$ is central, so it is fixed by $\Ad_{w}$ as well. The points $P_{2},P_{3},P_{4}$ are fixed by $\iota$ and permuted by $M_{0}$; the local form of $\iota$ at such a point is the inner automorphism by the surviving unit, and its fixed subalgebra is the corresponding complex line. The drive carries the lines with the points: $\Ad_{w}$ sends $i\mapsto j\mapsto k\mapsto i$ (certificate \texttt{L2\_Ad\_w\_cycles}, with $(2w)^{3}=-8$ giving order three, \texttt{L2\_w\_cubed}), so $\Ad_{w}\Ad_{i}\Ad_{w}^{-1}=\Ad_{j}$ and $\C_{i}\mapsto\C_{j}\mapsto\C_{k}$ exactly as $P_{2}\to P_{3}\to P_{4}$. A point with trivial stabiliser, which by Definition~\ref{def:A4}(iii) is a generic point, has nothing fixed and carries $\Hquat$.
\end{proof}

The reading in the statement, that the local form of $\iota$ is conjugation at the drive-fixed corner and the inner automorphism by the surviving unit at the three cycled corners, is the one place in this section where the proposition is \rigD\ rather than \rigR. It is the statement that the involution's action on the algebra varies by corner in the way the drive's three-cycle requires, and it is consistent with the corner lift of Corollary~\ref{cor:lift}, which distinguishes one corner by a sign; but the lift acts on spinors and the reading acts on the algebra, and the identification of the one with the other is a definition rather than a computation. We state it as such. The labelling of which cycled corner carries which unit is a convention and is fixed once in Appendix~\ref{app:labels}.

\begin{remark}[one thing, three shadows]\label{rem:three-shadows}
Proposition~\ref{prop:collapse} says more than that three corners carry three complex lines. The drive carries the line at one cycled corner to the line at the next, $\Ad_{w}\colon\C_{i}\to\C_{j}\to\C_{k}$, with the points: the three are one corner and its line, moved round a cycle of three. Seen from the place the drive leaves fixed (Definition~\ref{def:A5}(i)), that one thing is met at three places in each cycle and leaves three shadows. Under the fifth axiom the three are therefore one thing seen three times, and not three things. What those three shadows are in the physics, and why exactly three kinds of matter are observed, is a statement about matter, which this paper does not contain; here we record only that the geometry has one object where an observer counts three. \rigD\ (reading: Definition~\ref{def:A5})
\end{remark}

\subsection{Worked cases}\label{sec:fibre-worked}

\emph{The Hamilton product.} For $p=p_{1}i+p_{2}j+p_{3}k$ and $q=q_{1}i+q_{2}j+q_{3}k$ purely imaginary, $pq=-(\mathbf p\cdot\mathbf q)+\mathbf p\times\mathbf q$: the real part is minus the dot product, the imaginary part the cross product. The smallest case is $p=q=i$: $i^{2}=-1$, the whole product real. The imaginary part alone is not closed under multiplication, which is why the fibre algebra is all of $\Hquat=\R\oplus\ImH$ and not its imaginary part.

\emph{Compact and split corners.} On $\C_{u}=\R\oplus\R u$ the norm is $N(x+yu)=x^{2}-u^{2}y^{2}$. For $u=i$, $N=x^{2}+y^{2}$ and the unit group $N=1$ is the circle $U(1)$. In the split algebra $M_{2}(\R)$, with $j^{2}=+1$, $N=x^{2}-y^{2}$ on $\C_{j}$ and the unit group is the hyperbola $\{x^{2}-y^{2}=1\}$, non-compact; the element $1+j$ has $N=0$ and no inverse. A1 excludes the second picture at every corner, and with it every zero-divisor.

\emph{Fixed sets, by coordinates.} For $q=a+bi+cj+dk$: $\bar q=a-bi-cj-dk$, equal to $q$ iff $b=c=d=0$. $\Ad_{i}q=iqi^{-1}=-iqi=a+bi-cj-dk$, equal to $q$ iff $c=d=0$. $\overline{\Ad_{i}q}=a-bi+cj+dk$, equal to $q$ iff $b=0$: the fixed set $\{a+cj+dk\}$ contains $j$ and $k$ but $jk=i$, so it is not a subalgebra. Right-multiplication by $i$ composed with conjugation, $q\mapsto\overline{qi}$, fixes only the line $\R(1-i)$, which is not closed either ($(1-i)^{2}=-2i$). The four commutative fixed subalgebras are therefore exactly those of Proposition~\ref{prop:fixed}, and nothing else of the eight candidates qualifies.

\emph{The order-three automorphism.} With $2w=1+i+j+k$ and $\overline{2w}=1-i-j-k$, direct multiplication gives $(2w)\,i\,\overline{2w}=4j$, $(2w)\,j\,\overline{2w}=4k$, $(2w)\,k\,\overline{2w}=4i$, and $(2w)^{3}=-8$; dividing by $N(2w)=4$, $\Ad_{w}$ cycles $i\to j\to k\to i$ and $w^{3}=-1$. This is the algebraic image of the corner three-cycle of Proposition~\ref{prop:threecycle}, and it conjugates $\Ad_{i}$ to $\Ad_{j}$ to $\Ad_{k}$, so that the three complex corners are carried with the three cycled points.

\emph{Fixed points of inversion.} On $T^{n}$, $z\mapsto-z$ fixes the half-lattice $(\tfrac12\Z/\Z)^{n}$: two points on $S^{1}$, four on $\Ttwo$, eight on $T^{3}$. The Lefschetz number of $-I$ is $\det(I-(-I))=2^{n}$, in agreement. Four corners, the number forced by the algebra, select $n=2$.

\emph{The stabilisers.} Under $\langle\iota,M_{0}\rangle$ on $\Ttwo$: $P_{1}$ has stabiliser $\langle\iota,M_{0}\rangle$ and algebra $\Fix(\text{conj})\cap\Fix(\Ad_{w})=\R$; $P_{2},P_{3},P_{4}$ have stabiliser $\langle\iota\rangle$ and algebras $\C_{j},\C_{i},\C_{k}$ (Appendix~\ref{app:labels}); a point such as $(\tfrac13,0)$ is moved by $\iota$ to $(\tfrac23,0)$ and by $M_{0}$ to $(\tfrac23,\tfrac13)$, has trivial stabiliser, and carries $\Hquat$.

\subsection{The pillowcase and its cones}

\begin{theorem}[the fibre]\label{thm:fibre}
The quotient $\Ttwo/\iota$ is the pillowcase orbifold: a sphere with four cone points of cone angle $\pi$, the unique closed Euclidean two-orbifold of signature $(2,2,2,2)$. Its orbifold Euler characteristic is zero and its Euler characteristic is two. \rigR
\end{theorem}

\begin{proof}
$\iota$ has order two, so each fixed point is a cone point of order two, angle $2\pi/2=\pi$. A closed two-orbifold with four cone points of order two on a sphere has orbifold Euler characteristic $2-4\cdot\tfrac12=0$ (\texttt{L2\_chi\_orb\_zero}), hence is Euclidean \cite{Thurston}; among the seventeen Euclidean two-orbifolds the one with signature $(2,2,2,2)$ and no mirror boundary is unique, $\R^{2}/p2=\Ttwo/\Ztwo$. If the underlying surface were not a sphere the characteristic would be negative. For $\chi$, two routes: the quotient formula $\chi(X/\sigma)=\tfrac12(\chi(X)+\chi(\Fix\sigma))=\tfrac12(0+4)=2$ \cite{Bredon}; and orbifold Gauss--Bonnet with the flat metric, $2\pi\chi=\int K+\sum_{\alpha}\mathrm{deficit}_{\alpha}=0+4\pi$ \cite{Satake,Kawasaki1978}, so $\chi=2$ (\texttt{L2\_chi\_pillowcase}).
\end{proof}

\emph{The cone, by hand.} Near a fixed point $z\mapsto-z$ is rotation by $\pi$ about the point, so a small disc of angle $2\pi$ is folded to a cone of angle $\pi$ with deficit $\pi$. Four such deficits sum to $4\pi=2\pi\cdot2$, which is the Gauss--Bonnet count for $\chi=2$ with a flat interior. The alternative count, half the Euler characteristic of the torus plus half the number of fixed points, gives $\tfrac12(0+4)=2$ with no metric at all; that the two agree is the first consistency check between the metric and the combinatorial pictures of the fibre.

The four cone points are the geometry of Proposition~\ref{prop:collapse}: each is a point of angle deficit $\pi$ at which the fibre algebra has collapsed to a commutative line, the deficit being the trace of the identification $z\sim-z$ that produced the line. In the covering torus with periods $(L_{\psi},L_{\omega})$ any two distinct corners are at distance at least $\tfrac12\min(L_{\psi},L_{\omega})$: the corners are separated by the periods, and the ratio of the periods is the one continuous datum of the fibre, its aspect ratio $\tautwo=L_{\omega}/L_{\psi}$, whose fate is settled in Section~\ref{sec:theorem}.

\section{Dimension twelve}\label{sec:dimension}

\begin{theorem}[the dimension]\label{thm:dim}
Under Definition~\ref{def:levels} a geometry satisfying A1--A4 has dimension divisible by four and at least twelve; under criterion~(4) of Definition~\ref{def:minimal} its dimension is twelve. \rigD
\end{theorem}

\begin{proof}
An almost-quaternionic structure exists only in real dimension $4n$ \cite{Salamon1982}, so $4\mid D$. The time--space block $T\times\Sigma$ carries the quaternionic structure with the drive along $T$; its natural model is $\Hquat^{\times}=\R_{>0}\times S^{3}$, of dimension four, and this is the block's dimension (Section~\ref{sec:base} fixes $\Sigma$). Hence $\dim F=D-4$. A3 places the full algebra with its four corners on $F$ (Definition~\ref{def:levels}). On a four-manifold every oriented Riemannian structure is almost-quaternionic, since $\mathrm{Sp}(1)\mathrm{Sp}(1)=SO(4)$ and the triple $I,J,K$ is all of $\Lambda^{2}_{+}$: A3 would carry no information on $F$. So $F$ has quaternionic dimension at least two, $\dim F\geq8$, $D\geq12$ (\texttt{L3\_dim\_ge\_twelve}). Every $D$ from four to eleven fails: either $4\nmid D$, or $D\in\{4,8\}$ with $F$ of quaternionic dimension zero or one (\texttt{L3\_below\_twelve}); $D=12$ is admissible (\texttt{L3\_twelve\_ok}); $D=16,20,\dots$ are admissible and excluded by least dimension.
\end{proof}

\subsection{Worked cases}\label{sec:dim-worked}

The elimination, dimension by dimension. The time--space block has dimension four throughout, so $\dim F=D-4$ and the quaternionic dimension of $F$ is $m=(D-4)/4$ when this is an integer.

\begin{center}
\begin{tabular}{lllll}
\toprule
$D$ & $4\mid D$? & $\dim F$ & $m$ & verdict\\
\midrule
4 & yes & 0 & 0 & no internal factor at all\\
5, 6, 7 & no & -- & -- & no almost-quaternionic structure on $M$\\
8 & yes & 4 & 1 & $\mathrm{Sp}(1)\mathrm{Sp}(1)=SO(4)$: A3 vacuous on $F$\\
9, 10, 11 & no & -- & -- & no almost-quaternionic structure on $M$\\
12 & yes & 8 & 2 & admissible; least\\
13, 14, 15 & no & -- & -- & no almost-quaternionic structure on $M$\\
16 & yes & 12 & 3 & admissible; excluded by least dimension\\
\bottomrule
\end{tabular}
\end{center}

The case $D=8$ deserves a sentence, because non-commutativity alone does not exclude it: $\Hquat(2)$, the quaternionic $2\times2$ matrices that would act in quaternionic dimension two on the whole of $M$, is non-commutative, and an eight-manifold can carry an almost-quaternionic structure. What fails is the placement of the algebra on the internal factor. With four dimensions spent on the time--space block, the internal factor of an eight-dimensional geometry is a four-manifold, and on a four-manifold the three almost-complex structures $I,J,K$ span the self-dual two-forms of any oriented metric: the quaternionic structure is automatic, carries no information, and the four corners of Section~\ref{sec:fibre} have nowhere to live. Twelve is the first dimension at which the internal factor is large enough for A3 to say something about it.

The split within twelve: $\dim\Sigma+\dim(\text{base})=9$, $\dim\Sigma\in\{3,7\}$ from $\dim\Sigma+1\equiv0$, $\dim(\text{base})\in\{2,6\}$ from $\dim(\text{base})+2\equiv0$ (mod four). The pairs summing to nine are $(3,6)$ and $(7,2)$. The second makes $F$ four-dimensional and fails as $D=8$ did. So $\dim\Sigma=3$, $\dim(\text{base})=6$, $\dim F=8$.

\begin{remark}[the split]\label{rem:split}
Within $D=12$, $\dim\Sigma+\dim(\text{base})=9$ with $T\times\Sigma$ and $F=\text{base}\times\Ttwo/\Ztwo$ both quaternionic forces $(\dim\Sigma+1,\dim(\text{base})+2)\equiv(0,0)\pmod 4$. The congruences alone admit $(3,6)$ and $(7,2)$; the second gives $F$ four-dimensional, excluded by the same argument as $D=8$. With that condition the split is uniquely $(3,6)$ (\texttt{L3\_split\_three\_six}; the excluded alternative is recorded in \texttt{L3\_seven\_two\_meets\_congruences}). The certificate surfaced the second solution; the condition excluding it is stated rather than assumed.
\end{remark}

\section{The base}\label{sec:base}

\subsection{The sphere of imaginary units}

\begin{proposition}[the two-sphere]\label{prop:S2}
The elements $u\in\ImH$ with $u^{2}=-1$ form a two-sphere, on which the inner automorphisms $\Ad_{g}$, $g\in\mathrm{Sp}(1)$, act transitively through $\mathrm{Sp}(1)\to SO(3)$. If a $\Spinc$ structure selects one unit $i$, its stabiliser is $U(1)$ and the orbit of $i$ is $\mathrm{Sp}(1)/U(1)=\Stwo$. Under Definition~\ref{def:orbits} this sphere is a factor of the base. \rigD
\end{proposition}

\begin{proof}
$\ImH\cong\R^{3}$ and $u^{2}=-\lvert u\rvert^{2}$, so the condition is $\lvert u\rvert=1$. The three corner units are an orthonormal frame, the integer points of the sphere (\texttt{L4\_unit\_sphere\_axes}); $V$ acts on it as the rotations by $\pi$ about the three axes. The stabiliser of $i$ in $\mathrm{Sp}(1)$ is $\{e^{i\theta}\}$, and $\mathrm{Sp}(1)/U(1)$ is the Hopf quotient $S^{3}\to\Stwo$. The last sentence is Definition~\ref{def:orbits}.
\end{proof}

\subsection{The space of complex directions}

\begin{proposition}[the projective plane]\label{prop:CP2}
Let the three imaginary units be made complex-valued by the $\Spinc$ line of Proposition~\ref{prop:spinc} below, so that the corner data take values in $\C\otimes\ImH\cong\C^{3}$, and let the symmetry preserving the induced Hermitian form and the $\Spinc$ phase be $U(3)$. The space of complex lines in $\C^{3}$ is $\CPtwo=U(3)/(U(1)\times U(2))=SU(3)/U(2)$; its Euler characteristic three counts the units, its tangent bundle has $c_{1}=3h$, and its anticanonical line is $\Oline{3}$. Under Definition~\ref{def:orbits} it is a factor of the base. \rigD
\end{proposition}

\begin{proof}
The three coordinate lines $[1{:}0{:}0],[0{:}1{:}0],[0{:}0{:}1]$ are the units; the maximal torus of $U(3)$ fixes exactly these, and the Lefschetz count gives $\chi=3$ (\texttt{L4\_chi\_counts}). By the Euler sequence $0\to\mathcal{O}\to\mathcal{O}(1)^{\oplus3}\to T\CPtwo\to0$ and the Whitney formula, $c(T\CPtwo)=(1+h)^{3}=1+3h+3h^{2}$ in $\Z[h]/(h^{3})$ \cite{GriffithsHarris} (\texttt{L4\_chern\_CP2}); $\int c_{2}=3=\chi$ checks the smallest case. Two readings enter and are named: that complexification is forced, which is Proposition~\ref{prop:spinc}'s statement that spinor-valued data are complex; and that the geometry records directions rather than vectors, which is the scale-freeness of A3 together with the exemption of overall scale in Definition~\ref{def:minimal}. This proposition is the least direct of the section and is tagged accordingly.
\end{proof}

\subsection{The twist is necessary}

\begin{proposition}[$\Spinc$ is forced]\label{prop:spinc}
Under Definition~\ref{def:levels} a spinor bundle exists on $M$. Since $w_{2}(\CPtwo)=h\bmod2\neq0$, no spin structure exists on a geometry whose base contains $\CPtwo$, and spinors exist only through a $\Spinc$ structure, whose determinant line satisfies $c_{1}(L)\equiv w_{2}\pmod2$ and is therefore odd. $\Spinc$ is the least such structure. The structure induced by the almost-complex structure of A3's selected unit is the canonical one, with $L=K^{-1}=\Oline{3}$, and every other admissible line is $\Oline{3}\otimes E^{2}$ for a line $E$, excluded by criterion~(1) of Definition~\ref{def:minimal} as an unforced datum. \rigD
\end{proposition}

\begin{proof}
An oriented manifold carries spinors iff $w_{2}=0$ or $w_{2}$ lifts to an integral class, the latter being the $\Spinc$ condition with $c_{1}(L)\equiv w_{2}$ \cite{LawsonMichelsohn}. $w_{2}(\CPtwo)$ is the reduction of $c_{1}=3h$, which is $h\neq0$ (\texttt{L4\_not\_spin\_and\_congruence}); $3h$ is integral, so $\Spinc$ exists and $c_{1}(L)$ is odd. Among $\mathrm{Spin}\subset\Spinc\subset\mathrm{Spin}^{h}\subset\cdots$, $\Spinc$ has the least structure group lifting a non-zero $w_{2}$, and the larger ones reduce to it when $w_{2}$ lifts integrally. Every almost-complex manifold carries a canonical $\Spinc$ structure with determinant $K^{-1}$ \cite{LawsonMichelsohn,Friedrich}; on $\CPtwo$, $K^{-1}=\Oline{3}$; the $\Spinc$ structures form a torsor over $H^{2}(\Z)$ by $L\mapsto L\otimes E^{2}$, and $\Oline{1},\Oline{5}$ are $\Oline{3}\otimes\Oline{\mp1}^{2}$ (\texttt{L4\_odd\_lines}). The $\Stwo$ factor is spin and the pillowcase is flat, so neither obstructs. The flux parity, $Q_{c}$ odd on $\CPtwo$ and $Q_{w}$ even on $\Stwo$, is the same congruence on the flux side (\texttt{L4\_flux\_parity}).
\end{proof}

\begin{corollary}[the corner lift]\label{cor:lift}
With $c_{1}(L)$ odd the lift of $\iota$ to the spinor bundle acts at the corner $C_{m,n}=(\tfrac m2,\tfrac n2)$ by the sign $(-1)^{mn}$: three corners carry $+1$ and one carries $-1$, with product $-1$, the Atiyah--Bott fixed-point index $(-1)^{d}$ for odd determinant class \cite{AtiyahBott}. For even $L$ the product would be $+1$. \rigR
\end{corollary}

\begin{proof}
The eight functions $(\Z_{2})^{2}\to\Z_{2}$ vanishing at the origin split into the four characters, with product $+1$, and the four quadratic refinements of the intersection pairing, with product $-1$; the Atiyah--Bott input selects the refinement class for odd $d$. The sign table is \texttt{L4\_four\_corner\_product}.
\end{proof}

\subsection{Worked cases}\label{sec:base-worked}

\emph{The unit sphere and its axes.} $u=bi+cj+dk$ has $u^{2}=-(b^{2}+c^{2}+d^{2})$, so $u^{2}=-1$ is the unit sphere in $\R^{3}$. Its integer points are $\pm i,\pm j,\pm k$: the corner units, an orthonormal frame. $\Ad_{i}$ fixes $\pm i$ and sends $j\mapsto-j$, $k\mapsto-k$: rotation by $\pi$ about the $i$-axis; the Klein four-group $V$ is the group of such rotations about the three axes. The circle $U(1)=\{\cos\theta+i\sin\theta\}$ fixes $i$ and rotates the equator $\{cj+dk\}$, the units anticommuting with $i$, among which are $j$ and $k$.

\emph{The Euler sequence.} From the exact sequence
\[
0\to\mathcal{O}\to\mathcal{O}(1)^{\oplus3}\to T\CPtwo\to0
\]
and $c(\mathcal{O})=1$, $c(\mathcal{O}(1))=1+h$, the Whitney formula gives
\[
c(T\CPtwo)=(1+h)^{3}=1+3h+3h^{2}+h^{3},
\]
and $h^{3}=0$ on $\CPtwo$. So $c_{1}=3h$, $c_{2}=3h^{2}$, $\int c_{2}=3=\chi(\CPtwo)$, the Betti numbers being $1,0,1,0,1$. The anticanonical line is $\det T^{1,0}$ with $c_{1}=3h$, i.e.\ $\Oline{3}$; the canonical line is $\Oline{-3}$.

\emph{The admissible lines.} $c_{1}(L)\equiv w_{2}=h\pmod2$ admits $L=\Oline{n}$ for every odd $n$. The canonical structure has $L=K^{-1}=\Oline{3}$ and no further datum. Its neighbours are twists of it,
\[
\Oline{1}=\Oline{3}\otimes\Oline{-1}^{2},\qquad \Oline{5}=\Oline{3}\otimes\Oline{1}^{2},
\]
each requiring a twisting line $E=\Oline{\mp1}$: a discrete choice which criterion~(1) removes.

\emph{The eight corner functions.} A function $q:\{0,1\}^{2}\to\Z_{2}$ with $q(0,0)=0$ is one of eight. The four characters $q=am+bn$ have $\sum q=2(a+b)\equiv0$ and corner product $+1$; the four quadratic refinements $q=mn+am+bn$ have $\sum q\equiv1$ and product $-1$. For $q=mn$: values $0,0,0,1$ at $(0,0),(1,0),(0,1),(1,1)$, signs $+,+,+,-$. Three corners carry $+1$ and one carries $-1$, and it is the odd determinant class that puts the lift in the refinement class.

\emph{Flux parity.} The flux quantum through a closed surface $S$ is $Q_{S}=\int_{S}c_{1}(L)\equiv\langle w_{2},[S]\rangle$ mod two. Through the generating $\mathbb{CP}^{1}\subset\CPtwo$, $Q_{c}=3h\cdot[\mathbb{CP}^{1}]=3$, odd. Through $\Stwo$, which is spin, $Q_{w}$ is even, and the least even quantum compatible with one canonical unit $c_{1}(K^{-1}_{S^{2}})=2$ per sector is a multiple of two.

\subsection{The spatial factor}

\begin{proposition}[$\Sigma=S^{3}$]\label{prop:S3}
The time--space block $T\times\Sigma$ of Theorem~\ref{thm:dim} is modelled on $\Hquat^{\times}=\R_{>0}\times S^{3}$, so $\Sigma=S^{3}/\Gamma$ for a finite $\Gamma\subset\mathrm{Sp}(1)$; every such quotient is rigid, and the last clause of criterion~(4) selects $\Gamma$ trivial, $\Sigma=S^{3}=\mathrm{Sp}(1)$. \rigD
\end{proposition}

The quotients by $\Z_{2}$ and $\Z_{3}$ pass criteria (1)--(3) and are removed only by the least-$\pi_{1}$ clause, which is recorded as the reason; a stronger principle, that the geometry carries no unforced topology, would remove them without the clause and is not adopted.

\begin{theorem}[the base]\label{thm:base}
The base of a minimal geometry is $\CPtwo\times\Stwo$, its $\Spinc$ determinant line is $\Oline{3}$, its spatial factor is $S^{3}$, and the Euler characteristic of the internal factor is $\chi(\Keight)=\chi(\CPtwo)\,\chi(\Stwo)\,\chi(\Ttwo/\Ztwo)=3\cdot2\cdot2=12$. \rigD
\end{theorem}

\begin{proof}
Propositions~\ref{prop:S2}, \ref{prop:CP2}, \ref{prop:spinc}, \ref{prop:S3} supply the factors. The warping is a metric datum and does not change the homotopy type, so $\chi$ is multiplicative by K\"unneth \cite{Hatcher}, with $\chi(\Ttwo/\Ztwo)=2$ from Theorem~\ref{thm:fibre} (\texttt{L3\_chi\_K8}). The identity $12=1+3+8=\dim\Mtwelve$ holds and is recorded, not used.
\end{proof}

\section{The minimality theorem}\label{sec:theorem}

\subsection{Statement and proof}

The theorem collects the four sections into one statement. Its hypotheses are the five axioms, the six readings of Section~\ref{sec:axioms}, the minimality principle of Section~\ref{sec:minimal}, and one declared input; its conclusion names every factor and every discrete datum of the geometry. What a referee should check is that no factor is named before it is forced, and the proof is written so that each clause points to the section that forces it.

\begin{theorem}[minimality]\label{thm:main}
Among geometries $(T\times\Sigma)\times F$ satisfying A1--A5 under the readings of Definitions~\ref{def:A1}, \ref{def:A2}(i), \ref{def:levels}, \ref{def:orbits} and~\ref{def:A5}(i),~(ii), with the overall scale stabilised by a declared sector, the minimal geometry in the sense of Definition~\ref{def:minimal} is unique and is
\[
\Mtwelve=(S^{1}_{\Sol}\times S^{3})\times\bigl[(\CPtwo\times\Stwo)\times_{w}(\Ttwo/\Ztwo)\bigr]^{\Spinc},
\]
with monodromy $\bigl(\begin{smallmatrix}2&1\\1&1\end{smallmatrix}\bigr)$, determinant line $\Oline{3}$, $\chi(\Keight)=12$, dimension twelve, and no continuous modulus. The primes two, three and five of its topological integers are outputs of the selection.
\end{theorem}

\begin{proof}
A4 with criterion~(2) gives the drive (Theorem~\ref{thm:drive}). A1--A4 with criterion~(4) give the fibre, its algebra, its corners and the pillowcase (Section~\ref{sec:fibre}). A3 with criterion~(4) gives dimension twelve and the split $(3,6)$ (Theorem~\ref{thm:dim}, Remark~\ref{rem:split}). A3 with Definition~\ref{def:orbits} gives $\Stwo$ and $\CPtwo$, with $\Spinc$ necessary and $L=\Oline{3}$ by criterion~(1), and $\Sigma=S^{3}$ by criterion~(4) (Theorem~\ref{thm:base}). A5 gives the place we see from, $P_{1}$, the one point the drive returns to itself (Proposition~\ref{prop:threecycle} and the paragraph after it), and puts the growth law on the pillowcase alone (Definition~\ref{def:A5}(ii)). The moduli are treated in Proposition~\ref{prop:moduli}. The primes: two enters through the involution of A2 and $\chi(\Stwo)$; three through the three imaginary units, $\chi(\CPtwo)=3=c_{1}$; five through $\phi$ on the drive alone. Criterion~(3) therefore removes nothing that the other criteria had not removed, and is an output.
\end{proof}

The deflationary reading. The theorem asserts uniqueness among minimal realisations under the stated readings; it does not assert that a non-minimal geometry, one of dimension sixteen or with a lens-space spatial factor, fails the axioms, and it does not derive the overall scale. The readings are definitions, not results: a reader who takes Definition~\ref{def:orbits} in its fibrewise form, under which the sphere of units exists over each point but is not a factor, will not obtain Propositions~\ref{prop:S2} and~\ref{prop:CP2}, and the base is then determined among products only. Proposition~\ref{prop:CP2} is the least direct step and is tagged \rigD\ with its two readings named.

\subsection{The elimination table}

\begin{center}
\small
\begin{tabular}{@{}p{0.30\textwidth}p{0.33\textwidth}p{0.24\textwidth}@{\hspace{4pt}}l@{}}
\toprule
Competitor & Fails at & Where & Tag\\
\midrule
non-warped product & A2(i) & Definition~\ref{def:A2} & \rigD\\
$D\not\equiv0\ (4)$ & no quaternionic structure & Theorem~\ref{thm:dim} & \rigR\\
$D=4$, $D=8$ & $F$ of quaternionic dimension $<2$ & Theorem~\ref{thm:dim} & \rigR\\
$D=16,20,\dots$ & least dimension & Definition~\ref{def:minimal}(4) & \rigD\\
trace $2$ & entropy zero & Proposition~\ref{prop:hyp} & \rigR\\
trace $\geq4$ & least entropy; trace three unique & Theorem~\ref{thm:drive} & \rigR\\
$M_{2}(\R)$ as fibre algebra & non-compact corner unit groups & Proposition~\ref{prop:H} & \rigD\\
$T^{n}$, $n\neq2$, as fibre & $2^{n}\neq4$ corners & Proposition~\ref{prop:surface} & \rigR\\
$T^{4}/\Z_{2}$, K3 as internal factor & $16$ and $0$ corners; $D=8$ & Prop.~\ref{prop:corners}, Thm.~\ref{thm:dim} & \rigR\\
$S^{3}/\Z_{2}$, $S^{3}/\Z_{3}$ & least $\lvert\pi_{1}(\Sigma)\rvert$ & Proposition~\ref{prop:S3} & \rigD\\
spin bases: $S^{2}\times S^{2}\times S^{2}$, $\mathbb{CP}^{3}$, $S^{6}$ & $\Spinc$ unforced; no $\Stwo$ factor & Props.~\ref{prop:spinc}, \ref{prop:S2} & \rigD\\
bases with moduli: $\CPtwo\times T^{2}$, $T^{6}$ & criterion (1) & Definition~\ref{def:minimal} & \rigR\\
irreducible bases: the quadric $Q^{3}$ & no $\mathrm{Sp}(1)/U(1)$ factor; $\tautwo$ unfixed & Prop.~\ref{prop:S2}, Prop.~\ref{prop:moduli} & \rigD\\
$\Oline{1}$, $\Oline{5}$ & unforced twisting line & Proposition~\ref{prop:spinc} & \rigD\\
\bottomrule
\end{tabular}
\end{center}

\subsection{No continuous modulus}

\begin{proposition}[moduli]\label{prop:moduli}
Write $u=\tautwo^{2}=R_{w}^{2}/R_{c}^{2}$ for the square of the fibre aspect ratio and $x=R_{c}^{2}$ for the scale, with $R_{w},R_{c}$ the radii of $\Stwo$ and $\CPtwo$ tied by the warp to the two periods, and $Q_{w},Q_{c}$ the flux quanta through the two base factors. The stationarity conditions of the flux-plus-curvature energy are equivalent to $Q_{w}^{2}=2ux(6u+1)$ and $Q_{c}^{2}u=2x(18u+1)$; the scale is eliminated to the single constraint $Q_{w}^{2}(18u+1)=Q_{c}^{2}u^{2}(6u+1)$; with $(Q_{c},Q_{w})=(3,6)$, the least odd and the least even quanta compatible with Proposition~\ref{prop:spinc} and with one unit per corner-cycled sector, this is the integer cubic $6u^{3}+u^{2}-72u-4=0$, whose unique positive root fixes $\tautwo=1.8466$. No dimensionless modulus survives. The scale $x$ is fixed in the units of the energy but is a saddle of it, and is held by a sector the axioms do not supply; it is the declared input of Theorem~\ref{thm:main}. \rigD
\end{proposition}

\begin{proof}
The elimination of $x$ and the specialisation to the cubic are the identities \texttt{L5\_ratio} and \texttt{L5\_cubic}. An irreducible base has one radius and one flux quantum; the system then has a single equation after scale elimination and fixes nothing about $u$, so $\tautwo$ is a flat direction: this is the second route by which irreducible bases fail criterion~(1), independent of Proposition~\ref{prop:S2}. The stationarity conditions themselves are the Euler--Lagrange equations of the energy and enter as hypotheses here; their derivation belongs to the sequel.
\end{proof}

\subsection{Worked cases}\label{sec:theorem-worked}

\emph{The elimination of the scale.} The first stationarity condition is
\[
-12R_{w}^{2}-2R_{c}^{2}+Q_{w}^{2}R_{c}^{2}/R_{w}^{2}=0 ;
\]
multiply by $R_{w}^{2}$ and substitute $R_{w}^{2}=ux$, $R_{c}^{2}=x$. This gives $-12u^{2}x^{2}-2ux^{2}+Q_{w}^{2}x=0$, that is $Q_{w}^{2}=2ux(6u+1)$. The second, $-24R_{c}^{2}+Q_{c}^{2}-Q_{w}^{2}R_{c}^{4}/R_{w}^{4}=0$, becomes $Q_{c}^{2}=24x+Q_{w}^{2}/u^{2}=24x+2x(6u+1)/u=2x(18u+1)/u$. Dividing, $x$ cancels: $(Q_{w}/Q_{c})^{2}=u^{2}(6u+1)/(18u+1)$. At $(3,6)$ the left side is four and clearing denominators gives $6u^{3}+u^{2}-72u-4=0$.

\emph{The root.} $f(u)=6u^{3}+u^{2}-72u-4$ has $f(0)=-4$, $f(3)=-46$, $f(4)=108$; it is increasing for $u>2$ and has one positive root, bracketed by $f(3.409867)<0<f(3.409868)$. So $u^{*}=3.409868$, $\tautwo=\sqrt{u^{*}}=1.8466$, and $1/(1+u^{*})=0.22676$.

\emph{The flux table.} For $Q_{c}=3$ and the parity-admissible $Q_{w}\in\{2,4,6,8\}$ the constraint gives $u=1.103,\ 2.256,\ 3.410,\ 4.564$ respectively. Each is a zero-modulus geometry; the selection of $Q_{w}=6$ among them is one canonical unit per sector of the corner three-cycle, three units of $\chi(\Stwo)=2$, and is stated as such rather than argued from data.

\emph{The scale.} With $u^{*}$ fixed, $x=Q_{w}^{2}/(2u^{*}(6u^{*}+1))=36/(2\cdot3.4099\cdot21.459)\approx0.246$ in the units of the energy functional; this is a stationary value but, by the signature $(+,-)$ of the Hessian, a saddle along the scale direction. The scale is held by the declared sector and is the one measured input, the absolute pitch of a construction whose every ratio is fixed.

\subsection{Cross-check}

Witten's theorem \cite{Witten1981} states that the lowest-dimensional compact manifold admitting $SU(3)\times SU(2)\times U(1)$ as isometries is $\CPtwo\times\Stwo\times S^{1}=SU(3)/U(2)\times SU(2)/U(1)\times U(1)$. Propositions~\ref{prop:S2} and~\ref{prop:CP2} produce $\mathrm{Sp}(1)/U(1)$ and $SU(3)/U(2)$ from the three imaginary units with no reference to a gauge group. The two routes are logically independent and land on the same base; Witten's result is used only as this confirmation, never as a premise.

\section{Applications and outlook}\label{sec:outlook}

The geometry~\eqref{eq:M12} is now a theorem of four behavioural axioms and a minimality principle, and the derivation fixes more than the manifold. The four corners of the pillowcase are cone points of angle $\pi$, each the locus where the quaternionic algebra has collapsed to a commutative line, with the drive fixing one and cycling the other three (Propositions~\ref{prop:collapse}, \ref{prop:threecycle}); the determinant line is odd and the corner lift has product $-1$ (Corollary~\ref{cor:lift}); the aspect ratio is a root of an integer cubic (Proposition~\ref{prop:moduli}). These are the data that the papers of the programme consume in deriving the lepton and quark mixing structure from cone deficits and corner walks, and the present paper fixes their origin. We do not pursue those consequences here. \rigM

The sequel derives the field equation of the programme, with its coupling to the Coherent Field, from the same axioms on the geometry derived here; the two sheets of the involution of A2 and the non-commutative structure of A3 are the variables in which that equation is written.

\section*{Acknowledgements}
The author used Claude (Anthropic) and Gemini (Google DeepMind) as mathematical development and analytical tools during the preparation of this work. All intellectual content, theoretical framework, and scientific claims are the sole responsibility of the author.

\appendix

\section{Machine-verified certificates}\label{app:lean}

Five files, core library only, no Mathlib, no author axioms, zero \texttt{sorry}; fifty theorems. The files were compiled independently on two toolchains, the core toolchain used in drafting and the author's Lean project, with identical results; no proof depends on the default simplifier set of a particular Lean version. Classical theorems (Franks, Frobenius, Thurston, Latimer--MacDuffee, Lawson--Michelsohn, Salamon, Atiyah--Bott) are named in file headers and are never Lean axioms; the certificates establish the finite and arithmetic content those theorems consume.

\texttt{M1\_L1\_drive.lean} (Section~\ref{sec:drive}): $M_{0}\in\mathrm{SL}_{2}(\Z)$ of trace three; $SM_{0}S^{-1}=M_{0}^{-1}$; $t^{2}>4$, $t>0\Rightarrow t\geq3$; reduced forms of discriminant five are $(\pm1,1,\mp1)$, one cycle; odd trace $\Leftrightarrow$ order three in $\mathrm{SL}_{2}(\mathbb{F}_{2})$, by finite check; $\bar M_{0}$ of order three; $M_{0}=F^{2}$ with $F$ the Fibonacci matrix, the Fibonacci recursion in the entries, and $F\notin\mathrm{SL}_{2}(\Z)$; $\det(M_{0}-I)=-1$, so the drive returns one point only to itself.

\texttt{M1\_L2\_fibre.lean} (Section~\ref{sec:fibre}): Hamilton's relations; $ij\neq ji$; $N(q)=q\bar q$, positive-definite; fixed sets of conjugation and of $\Ad_{i},\Ad_{j},\Ad_{k}$; $\iota\circ\Ad_{i}$ fixes $i^{\perp}$, not closed; $\Ad_{w}$ cycles $i\to j\to k$, $(2w)^{3}=-8$; corner count four; $2^{n}=4\Rightarrow n=2$; $\chi_{\mathrm{orb}}=0$; $\chi=2$.

\texttt{M1\_L3\_dimension.lean} (Section~\ref{sec:dimension}): $D=4+4m$, $m\geq2\Rightarrow D\geq12$; dimensions four to eleven fail; twelve admissible; split $(3,6)$ unique given $\dim F\geq8$, and $(7,2)$ recorded as the congruence solution excluded; $3\cdot2\cdot2=12=1+3+8$.

\texttt{M1\_L4\_base.lean} (Section~\ref{sec:base}): integer points of the unit sphere of $\ImH$ are the three axes; $(1+h)^{3}$ coefficients; $3\bmod2=1$ and the congruence; odd lines $3+2e$; flux parity; four-corner product $-1$; $\chi(\CPtwo)=3$ as a count.

\texttt{M1\_L5\_vacuum.lean} (Section~\ref{sec:theorem}): elimination of the scale to $Q_{w}^{2}(18u+1)=Q_{c}^{2}u^{2}(6u+1)$; specialisation at $(3,6)$ to $6u^{3}+u^{2}-72u-4=0$.

The five files follow verbatim.

\subsection*{\texttt{M1\_L1\_drive.lean} --- the drive}
\lstinputlisting{M1_lean/M1_L1_drive.lean}

\subsection*{\texttt{M1\_L2\_fibre.lean} --- the fibre}
\lstinputlisting{M1_lean/M1_L2_fibre.lean}

\subsection*{\texttt{M1\_L3\_dimension.lean} --- the dimension}
\lstinputlisting{M1_lean/M1_L3_dimension.lean}

\subsection*{\texttt{M1\_L4\_base.lean} --- the base}
\lstinputlisting{M1_lean/M1_L4_base.lean}

\subsection*{\texttt{M1\_L5\_vacuum.lean} --- the moduli}
\lstinputlisting{M1_lean/M1_L5_vacuum.lean}

\section{The labelling of the corners}\label{app:labels}

The four corners are named by position, $P_{1}=(0,0)$, $P_{2}=(\tfrac12,0)$, $P_{3}=(0,\tfrac12)$, $P_{4}=(\tfrac12,\tfrac12)$ in coordinates $(\theta_{\psi},\theta_{\omega})$, and by algebra, $C_{0}=\R$, $C_{i}$, $C_{j}$, $C_{k}$. Throughout the programme the dictionary is: $P_{1}=C_{0}$ (fixed by the drive; $\iota$ acts as conjugation); $P_{3}=C_{i}$ (the $\omega$ half-period); $P_{2}=C_{j}$ (the $\psi$ half-period); $P_{4}=C_{k}$ (the corner at which the $\Spinc$ lift carries the sign $-1$). The drive fixes $P_{1}$; the lift distinguishes $P_{4}$; these are different structures and distinguish different corners, and the drive permutes the lift's sectors. The assignment of $j$ and $k$ between $P_{2}$ and $P_{3}$ is the orientation of the three-cycle, which Theorem~\ref{thm:drive} shows to be a single conjugacy class and which no equivariant invariant can detect: it is a convention and carries no content.

\begin{thebibliography}{99}
\bibitem{AtiyahBott} M.~F. Atiyah and R. Bott, A Lefschetz fixed point formula for elliptic complexes II, \emph{Ann. of Math.} 88 (1968) 451--491.
\bibitem{Bredon} G.~E. Bredon, \emph{Introduction to Compact Transformation Groups}, Academic Press, 1972.
\bibitem{Franks1970} J. Franks, Anosov diffeomorphisms, \emph{Proc. Sympos. Pure Math.} 14 (1970) 61--93.
\bibitem{FranksNewhouse} J. Franks, Anosov diffeomorphisms on tori, \emph{Trans. Amer. Math. Soc.} 145 (1969) 117--124; S. Newhouse, On codimension one Anosov diffeomorphisms, \emph{Amer. J. Math.} 92 (1970) 761--770.
\bibitem{Friedrich} T. Friedrich, \emph{Dirac Operators in Riemannian Geometry}, AMS, 2000.
\bibitem{Frobenius1878} G. Frobenius, \"Uber lineare Substitutionen und bilineare Formen, \emph{J. Reine Angew. Math.} 84 (1878) 1--63.
\bibitem{Gauss1801} C.~F. Gauss, \emph{Disquisitiones Arithmeticae}, 1801, \S\S183--205.
\bibitem{GriffithsHarris} P. Griffiths and J. Harris, \emph{Principles of Algebraic Geometry}, Wiley, 1978.
\bibitem{Hatcher} A. Hatcher, \emph{Algebraic Topology}, Cambridge University Press, 2002.
\bibitem{KatokHasselblatt} A. Katok and B. Hasselblatt, \emph{Introduction to the Modern Theory of Dynamical Systems}, Cambridge University Press, 1995.
\bibitem{Kawasaki1978} T. Kawasaki, The signature theorem for V-manifolds, \emph{Topology} 17 (1978) 75--83.
\bibitem{LatimerMacDuffee1933} C.~G. Latimer and C.~C. MacDuffee, A correspondence between classes of ideals and classes of matrices, \emph{Ann. of Math.} 34 (1933) 313--316.
\bibitem{LawsonMichelsohn} H.~B. Lawson and M.-L. Michelsohn, \emph{Spin Geometry}, Princeton University Press, 1989.
\bibitem{LyusternikFet} L. Lyusternik and A. Fet, Variational problems on closed manifolds, \emph{Dokl. Akad. Nauk SSSR} 81 (1951) 17--18.
\bibitem{Robertson1929} H.~P. Robertson, The uncertainty principle, \emph{Phys. Rev.} 34 (1929) 163--164.
\bibitem{Salamon1982} S. Salamon, Quaternionic K\"ahler manifolds, \emph{Invent. Math.} 67 (1982) 143--171.
\bibitem{Sarnak1982} P. Sarnak, Class numbers of indefinite binary quadratic forms, \emph{J. Number Theory} 15 (1982) 229--247.
\bibitem{Satake} I. Satake, The Gauss--Bonnet theorem for V-manifolds, \emph{J. Math. Soc. Japan} 9 (1957) 464--492.
\bibitem{Taussky1949} O. Taussky, On a theorem of Latimer and MacDuffee, \emph{Canad. J. Math.} 1 (1949) 300--302.
\bibitem{Thurston} W.~P. Thurston, \emph{The Geometry and Topology of Three-Manifolds}, Princeton lecture notes, 1980, Chapter 13.
\bibitem{Witten1981} E. Witten, Search for a realistic Kaluza--Klein theory, \emph{Nucl. Phys. B} 186 (1981) 412--428.
\end{thebibliography}

\end{document}
